\documentclass{article}
\usepackage[utf8]{inputenc}
\usepackage[russian]{babel}
\usepackage{amsmath}
\usepackage{amssymb}
\usepackage{mathtools}
\usepackage[a4paper, left=1.5cm, right=1.5cm, top=1.5cm, bottom=1.5cm]{geometry}\setlength{\parindent}{0pt}

\title{ Билет 8. Переход к пределу под знаком интеграла. Условия непрерывности интеграла, зависящего от параметра.}
\author{}
\date{}

\begin{document}

\maketitle

\section*{Переход к пределу под знаком интеграла}

Пусть $f(x,y)$ определена на $[a,b] \times Y$, где $y_0$ — предельная точка множества $Y$ (конечная/бесконечная, правая/левая/двусторонняя).

Предположим, что $\forall y \in Y$ существует интеграл $\int_a^b f(x,y) \, dx$.

Также пусть при $x \in [a,b]$ выполнено:
\[
f(x,y) \xrightarrow[y \to y_0]{} \varphi(x)
\]

Определим функцию:
\[
G(y) = \int_a^b f(x,y) \, dx
\]

\textbf{Требуется доказать:}
\[
\lim_{y \to y_0} G(y) = \int_a^b \varphi(x) \, dx
\]
или, что то же самое:
\[
\lim_{y \to y_0} \int_a^b f(x,y) \, dx = \int_a^b \lim_{y \to y_0} f(x,y) \, dx
\]

\subsection*{Доказательство}

Т.к. $\exists \int_a^b f(x,y) \, dx$ (по условию), то $\int_a^b \varphi(x) \, dx$ существует.

\textbf{Определение предела:} $\lim_{y \to y_0} G(y) = \int_a^b \varphi(x) \, dx$ означает:

$\forall \varepsilon > 0 \quad \exists \delta > 0 : \forall y \in Y \text{ таких, что } 0 < |y - y_0| < \delta$,

выполняется:
\[
\left| \int_a^b f(x,y) \, dx - \int_a^b \varphi(x) \, dx \right| < \varepsilon
\]

Используем неравенство:
\[
\left| \int_a^b f(x,y) \, dx - \int_a^b \varphi(x) \, dx \right| \leq \int_a^b |f(x,y) - \varphi(x)| \, dx
\]

Выберем $\frac{\varepsilon}{b-a} > 0$. Тогда $\exists \delta > 0 : \forall y \in Y$, при $0 < |y - y_0| < \delta$ и $\forall x \in [a,b]$ выполняется:
\[
|f(x,y) - \varphi(x)| < \frac{\varepsilon}{b-a}
\]

Умножаем на длину отрезка $(b-a)$:
\[
\int_a^b |f(x,y) - \varphi(x)| \, dx < \int_a^b \frac{\varepsilon}{b-a} \, dx = \varepsilon
\]

Следовательно,
\[
\left| \int_a^b f(x,y) \, dx - \int_a^b \varphi(x) \, dx \right| < \varepsilon
\]

Что и требовалось доказать.

\section*{Теорема о непрерывности интеграла по параметру}

Пусть $f(x,y)$ определена на $[a,b] \times [c,d]$.

Пусть $f(x,y)$ непрерывна на этом декартовом произведении (по двум переменным).

Тогда функция
\[
G(y) = \int_a^b f(x,y) \, dx
\]
является непрерывной на промежутке $[c,d]$.

\subsection*{Доказательство}

Если функция непрерывна по двум переменным, то она непрерывна как по $x$, так и по $y$, причём для каждого фиксированного $y$ она интегрируема по $x$.

Нам нужно показать, что $\forall y_0 \in [c,d]$:
\[
\lim_{y \to y_0} G(y) = G(y_0) = \int_a^b f(x,y_0) \, dx
\]

Это верно, но чтобы это доказать, нам нужно показать, что:
\[
\forall x \in [a,b] \quad f(x,y) \xrightarrow[y \to y_0]{} f(x,y_0)
\]

Так как $[a,b] \times [c,d]$ — замкнутое множество, и $f$ непрерывна на нём, то по теореме Кантора (о равномерной непрерывности на компакте):

$\forall \varepsilon > 0 \quad \exists \delta > 0 : \forall x', x'' \in [a,b], \forall y', y'' \in [c,d]$, таких что $|x' - x''| < \delta$ и $|y' - y''| < \delta$, выполняется:
\[
|f(x',y') - f(x'',y'')| < \varepsilon
\]

В частности, если положить $x' = x'' = x$, то получим:

$\forall \varepsilon > 0 \quad \exists \delta > 0 : \forall x \in [a,b], \forall y', y'' \in [c,d]$ таких, что $|y' - y''| < \delta$, выполняется:
\[
|f(x,y') - f(x,y'')| < \varepsilon
\]

В частном случае, если $y' = y_0$, то:

$\forall \varepsilon > 0 \quad \exists \delta > 0 : \forall x \in [a,b], \forall y'' \in [c,d]$ таких, что $|y_0 - y''| < \delta$, выполняется:
\[
|f(x,y_0) - f(x,y'')| < \varepsilon
\]

Или, что то же самое (меняя обозначения):

$\forall \varepsilon > 0 \quad \exists \delta > 0 : \forall x \in [a,b], \forall y \in [c,d]$ таких, что $|y_0 - y| < \delta$, выполняется:
\[
|f(x,y_0) - f(x,y)| < \varepsilon
\]

То есть, $f(x,y) \xrightarrow[y \to y_0]{} f(x,y_0)$ равномерно по $x \in [a,b]$.

Тогда, как показано выше, можно перейти к пределу под знаком интеграла:

\[
\lim_{y \to y_0} \int_a^b f(x,y) \, dx = \int_a^b \lim_{y \to y_0} f(x,y) \, dx = \int_a^b f(x,y_0) \, dx = G(y_0)
\]

Следовательно, $G(y)$ непрерывна в точке $y_0$. Так как $y_0$ — произвольная точка из $[c,d]$, то $G(y)$ непрерывна на всём отрезке $[c,d]$.

\end{document}
